% ============================================================================
% MAC F244 project report: Pairs trading with mean-reverting stochastic processes
% TEMPORARY MASTER COPY. Each section is tagged with its owner (M1-M5).
% Edit only your own sections. Red [TODO] notes mark what is still missing.
% Page budget: 4 pages of content (references outside the budget unless the
% instructor says otherwise). Numbers come from results/ written by notebook.ipynb.
% ============================================================================
\documentclass[11pt]{article}
\usepackage[margin=2cm]{geometry}
\usepackage{amsmath,amssymb}
\usepackage{graphicx}
\usepackage{booktabs}
\usepackage{xcolor}
\usepackage[authoryear,round]{natbib}
\usepackage[colorlinks=true,linkcolor=blue,citecolor=blue,urlcolor=blue]{hyperref}

\graphicspath{{../figs/}}
\newcommand{\todo}[1]{\textcolor{red}{[TODO: #1]}}
% \fig{name}{width}: shows the figure once notebook.ipynb has written figs/name.eps
\newcommand{\fig}[2]{\IfFileExists{../figs/#1.eps}{\includegraphics[width=#2\linewidth]{#1}}{\fbox{\parbox{#2\linewidth}{\centering\vspace{1.5cm}figure #1 pending\vspace{1.5cm}}}}}
\newcommand{\erfi}{\operatorname{erfi}}

\title{Pairs Trading with Mean-Reverting Stochastic Processes:\\
Ornstein--Uhlenbeck Modelling, ADF-Based Cointegration Selection\\ and Optimal Trading Bands}
\author{\todo{Name (ID)}, \todo{Name (ID)}, \todo{Name (ID)}, \todo{Name (ID)}, \todo{Name (ID)}\\
MAC F244, Semester I, 2026--27}
\date{\todo{submission date}}

\begin{document}
\maketitle

% ----------------------------------------------------------------------------
% OWNER: M5 (about 1/3 page)
\section{Introduction}
Pairs trading bets that the price gap between two related stocks returns to its
mean. It works only if that gap really is mean-reverting and if it is traded at the
right levels. We model the log-price spread of a pair as an Ornstein--Uhlenbeck (OU)
process, select pairs with augmented Dickey--Fuller (ADF) and Engle--Granger tests,
and compare four published trading rules on US and Indian stocks from 2010 to 2026.

We ask three questions: (i) how reliably can ADF and Engle--Granger tests find
mean-reverting spreads, (ii) how well does an OU model describe them and how
accurately can it be estimated, and (iii) do OU-based rules beat a simple distance
rule after transaction costs, across markets and periods?

\todo{M5: literature paragraph, four sentences, from \citet{gatev2006}, \citet{krauss2017},
\citet{rad2016} and \citet{dofaff2010} (task T07).}

% ----------------------------------------------------------------------------
% OWNER: M3 (OU model, ADF, Engle-Granger); M4 (trading rules, Bertram). About 1 page.
\section{Mathematical model}

\paragraph{Spread.} For a pair $(A,B)$ with log prices $a_t=\ln P^A_t$ and
$b_t=\ln P^B_t$, the spread is $X_t = a_t - \beta b_t - \mu$, where $\beta$ and $\mu$
come from an OLS regression on the formation window only.

\paragraph{Ornstein--Uhlenbeck process.} We model
\begin{equation}
  dX_t = \kappa(\theta - X_t)\,dt + \sigma\,dW_t, \qquad \kappa>0
  \label{eq:ou}
\end{equation}
\citep{uhlenbeck1930}. Applying It\^o's lemma to $e^{\kappa t}X_t$ removes the drift:
\begin{equation}
  X_t = \theta + (X_0-\theta)e^{-\kappa t} + \sigma\int_0^t e^{-\kappa(t-s)}\,dW_s ,
\end{equation}
so $X_t$ is Gaussian with mean $\theta+(X_0-\theta)e^{-\kappa t}$ and variance
$\sigma^2(1-e^{-2\kappa t})/(2\kappa)$. The stationary law is $N(\theta,\sigma^2/2\kappa)$
and a deviation halves in $\ln 2/\kappa$.

\paragraph{Exact discretisation and estimation.} With daily step $\Delta=1/252$,
\begin{equation}
  X_{t+\Delta} = a + bX_t + \varepsilon_t,\quad b=e^{-\kappa\Delta},\quad a=\theta(1-b),\quad
  \varepsilon_t\sim N\!\Big(0,\ \sigma^2\tfrac{1-e^{-2\kappa\Delta}}{2\kappa}\Big).
  \label{eq:ar1}
\end{equation}
OLS of $X_{t+\Delta}$ on $X_t$ is therefore the conditional maximum-likelihood
estimator, and
\begin{equation}
  \hat\kappa=-\frac{\ln\hat b}{\Delta},\qquad \hat\theta=\frac{\hat a}{1-\hat b},\qquad
  \hat\sigma=\hat s_\varepsilon\sqrt{\frac{2\hat\kappa}{1-\hat b^{2}}},\qquad
  \sigma_{eq}=\frac{\hat s_\varepsilon}{\sqrt{1-\hat b^{2}}},
\end{equation}
valid only when $0<\hat b<1$. In short samples $\hat\kappa$ is biased upward
\citep{yu2012}; Section~\ref{sec:results} measures the bias for our window lengths.
Euler--Maruyama gives $b\approx 1-\kappa\Delta$; because the noise in \eqref{eq:ou} is
additive, Milstein's correction vanishes and Milstein coincides with Euler--Maruyama
\citep{kloeden1992}.

\paragraph{ADF is a test of $\kappa=0$.} Subtracting $X_{t-1}$ from \eqref{eq:ar1}
gives the ADF regression \citep{dickeyfuller1979,saiddickey1984}
\begin{equation}
  \Delta X_t = \alpha + \gamma X_{t-1} + \sum_{i=1}^{p}\delta_i\,\Delta X_{t-i} + \varepsilon_t,
  \qquad \gamma=b-1=e^{-\kappa\Delta}-1 .
\end{equation}
$H_0:\gamma=0$ is exactly $H_0:\kappa=0$ (a random walk, no mean reversion), against
$H_1:\gamma<0$, i.e.\ $\kappa>0$. The $t$-statistic of $\gamma$ follows the
Dickey--Fuller distribution; the lag order $p$ is chosen by AIC.

\paragraph{Engle--Granger cointegration.} Both $a_t$ and $b_t$ are $I(1)$; we regress
$a_t=\mu+\beta b_t+u_t$ and test $\hat u_t$ for a unit root \citep{englegranger1987}.
Because $\hat\beta$ is fitted to make $\hat u_t$ look stationary, plain ADF critical
values over-reject (the spurious regression of \citealp{grangernewbold1974}); we use
Engle--Granger critical values \citep{mackinnon2010}. The test depends on which stock
is the regressand, so a pair is kept only if both directions reject at 5\%.

% OWNER: M4 from here to the end of Section 2
\paragraph{Trading rules.} The s-score of \citet{avellaneda2010} is
$s_t=(X_t-\hat\theta)/\sigma_{eq}$: open long when $s<-1.25$, open short when $s>1.25$,
close a long when $s>-0.50$ and a short when $s<0.75$. For the optimal bands of
\citet{bertram2010}, take the centred process $dX=-\kappa X\,dt+\sigma\,dW$, entry $a$,
exit $m>a$ and round-trip cost $c$. Solving
$\tfrac12\sigma^2u''-\kappa x u'=-1$ for first-passage times gives the expected cycle
length and the expected return per unit time
\begin{equation}
  \mathbb{E}[T]=\frac{\pi}{\kappa}\Big[\erfi\Big(\frac{m\sqrt\kappa}{\sigma}\Big)-\erfi\Big(\frac{a\sqrt\kappa}{\sigma}\Big)\Big],
  \qquad \mu(a,m)=\frac{m-a-c}{\mathbb{E}[T]} .
\end{equation}
With symmetric bands $a=-m$ we maximise $\mu$ over $m$ numerically.
\todo{M4: one sentence on the distance rule of \citet{gatev2006} and the z-score rule.}

\paragraph{Position P\&L.} Holding one unit of spread means \$1 long $A$ and \$$\beta$
short $B$; the daily return on gross capital is $(r^A_t-\beta r^B_t)/(1+|\beta|)$,
minus costs on every trade.

% ----------------------------------------------------------------------------
% OWNER: M2 (data), M1 (protocol). About 1/2 page.
\section{Data and methodology}
\paragraph{Data.} Daily adjusted closes from Yahoo Finance \citep{yahoo2026} for 56 US large caps in 7
sectors (Jan 2010--Sep 2026) and 26 NSE stocks in 5 sectors (Jan 2012--Sep 2026).
Pairs are formed only within a sector and a market (196 US candidate pairs per window).
The sample uses today's large caps, so it has survivorship bias.

\paragraph{Walk-forward protocol.} Pairs are selected on a 252-day formation window
and traded for the next 126 days with frozen parameters; the window then rolls by
126 days \citep{gatev2006}. Selection keeps Engle--Granger $p<0.05$ in both
directions, $\beta>0$, a half-life of 2--63 days and at least 12 mean crossings a
year \citep{sarmento2020}, ranked by the Engle--Granger statistic, top $K=20$.
Signals at the close of day $t$ are executed at the close of $t+1$, open positions are
closed at the end of the window, and costs are 10\,bp (US) or 15\,bp (India) per leg
per transaction. Thresholds are taken from the papers and never tuned; design choices
were frozen on 2010--2021 and 2022--2026 is reported as an untouched test block.

\begin{table}[h]
\centering\small
\begin{tabular}{llll}
\toprule
ID & Strategy & Selection & Entry and exit \\
\midrule
S0 & Distance \citep{gatev2006} & Minimum SSD & Open at $2\sigma$, close at crossing \\
S1 & Cointegration z-score & EG + ADF + OU filters & Open $|z|>2$, close $|z|<0.5$ \\
S2 & OU s-score \citep{avellaneda2010} & as S1 & $1.25 / 0.50 / 0.75$ \\
S3 & OU optimal bands \citep{bertram2010} & as S1 & Maximise $\mu(-m,m)$ at cost $c$ \\
\bottomrule
\end{tabular}
\caption{Strategies compared under one protocol.}
\end{table}

\paragraph{Evaluation.} Annualised return and volatility, Sharpe ratio with the
standard error of \citet{lo2002}, Newey--West $t$-statistic of the mean daily return
\citep{neweywest1987}, maximum drawdown, hit rate and market beta.

% ----------------------------------------------------------------------------
% OWNER: M5 (editor) with M3 (simulation), M1/M4 (backtest), M2 (persistence). About 1 2/3 pages.
\section{Results}\label{sec:results}

\paragraph{Simulation study.} \todo{M3: E1 bias of $\hat\kappa$ at 252 days; E2: which
half-lives a 12-month window detects, and the actual size of the naive residual test
versus Engle--Granger (from results/e2\_size.csv).}

\begin{figure}[h]
\centering
\fig{fig_e1_ou_estimation}{0.48}\hfill\fig{fig_e2_adf_power}{0.48}
\caption{Left: relative bias of $\hat\kappa$ against window length \todo{finding}.
Right: ADF rejection rate over half-life and window length \todo{finding}.}
\end{figure}

\paragraph{Headline backtest.} \todo{M4/M1: Table 1 and Figure 2, which rule wins and
whether the Sharpe ratio is significant; market beta.}

\begin{table}[h]
\centering\small
\begin{tabular}{lrrrrrr}
\toprule
Strategy & Ann.\ ret.\ \% & Vol \% & Sharpe (SE) & NW $t$ & Max DD \% & Beta \\
\midrule
S0 Distance & \todo{} & & & & & \\
S1 z-score & \todo{} & & & & & \\
S2 s-score & \todo{} & & & & & \\
S3 Bertram & \todo{} & & & & & \\
\bottomrule
\end{tabular}
\caption{US, 2010--2026, 10\,bp per leg (results/table1\_us.csv).}
\end{table}

\begin{figure}[h]
\centering
\fig{fig_e5_cumulative}{0.62}\hfill\fig{fig_e6_costs}{0.36}
\caption{Left: cumulative return on committed capital, test block shaded \todo{finding}.
Right: Sharpe ratio against cost per leg \todo{break-even costs}.}
\end{figure}

\paragraph{Robustness.} \todo{M5: E7 sensitivity, E8 India / sectors / sub-periods
(Table 2), M2: E9 share of pairs still cointegrated out of sample.}

% ----------------------------------------------------------------------------
% OWNER: M5. About 1/3 page.
\section{Conclusion}
\todo{M5: answer the three questions with numbers; limitations: survivorship bias,
fixed $\beta$ within a window, daily closes; future work: Kalman hedge ratio,
\citet{leungli2015} optimal stopping.}

\bibliographystyle{plainnat}
\bibliography{refs}
\end{document}
